% !TEX root = ../main.tex
\section{Preguntas de análisis}

A continuación presentamos las respuestas técnicas y directas a las doce preguntas de análisis formuladas en la guía de laboratorio:

\smallskip
\noindent\textbf{1. ¿Qué información pertenece al árbol representado y cuál a los árboles auxiliares?}\\
El \textbf{árbol representado} almacena la estructura lógica del grafo (las aristas reales que unen vértices y determinan componentes conexas). Los \textbf{árboles auxiliares (Splay Trees)} mantienen únicamente la partición en caminos preferidos; su recorrido \textit{in-order} codifica la profundidad relativa en el camino y sus nodos conservan marcas diferidas (\texttt{reversed}) y agregados asociativos (\texttt{path\_sum}, \texttt{size}).

\smallskip
\noindent\textbf{2. ¿Por qué una rotación splay no modifica las aristas reales del bosque?}\\
Porque las rotaciones se ejecutan exclusivamente dentro de los árboles auxiliares reconfigurando punteros de búsqueda binaria locales. Dado que toda rotación conserva con exactitud el orden simétrico (\textit{in-order}), la secuencia de vértices que integran el camino lógico en el árbol representado permanece inalterada.

\smallskip
\noindent\textbf{3. ¿Qué resultado estructural deja \texttt{EXPOSE(v)}?}\\
Agrupa todas las aristas desde la raíz del árbol representado hasta $v$ como aristas preferidas dentro de un \textbf{único Splay Tree auxiliar}, y posiciona al vértice $v$ en la \textbf{raíz de dicho auxiliar}, asegurando que no tenga hijo derecho hacia descendientes ($v.\text{right} = \text{None}$).

\smallskip
\noindent\textbf{4. ¿Por qué \texttt{MAKE-ROOT(v)} requiere invertir el camino expuesto?}\\
En un árbol dirigido hacia la raíz, cada nodo apunta hacia sus ancestros. Al transformar a $v$ en la nueva raíz de su componente, la dirección de todas las aristas en el camino raíz $\to v$ debe invertirse. Como el recorrido \textit{in-order} del árbol auxiliar mapea dicho camino de menor a mayor profundidad, conmutar la orientación mediante \texttt{reversed = True} efectúa la inversión completa en tiempo $O(1)$.

\smallskip
\noindent\textbf{5. ¿Qué errores aparecen si no se propaga \texttt{reversed} antes de rotar?}\\
Si se rota un nodo con la inversión pendiente, la rotación reordena los punteros antes de intercambiar los hijos, destruyendo irreversiblemente la propiedad \textit{in-order}. Esto produce que \texttt{path\_sum} acumule valores de vértices ajenos al camino, que \texttt{find\_root} descienda hacia hojas arbitrarias y que \texttt{cut} desconecte aristas incorrectas.

\smallskip
\noindent\textbf{6. ¿Cómo evita \texttt{LINK} la formación de ciclos?}\\
Antes de crear cualquier enlace, invoca \texttt{connected(u, v)}, la cual determina las raíces representadas mediante \texttt{find\_root}. Si $u$ y $v$ comparten la misma raíz o son idénticos, pertenecen ya a la misma componente conexa; la operación se aborta de inmediato lanzando una excepción \texttt{ValueError}, preservando la aciclicidad del bosque.

\pagebreak
\noindent\textbf{7. ¿Cómo comprueba \texttt{CUT} que dos nodos forman una arista directa?}\\
Ejecuta \texttt{make\_root(u)} y posteriormente \texttt{expose(v)}. Si existe una arista física entre ambos, $u$ debe ser el predecesor inmediato de $v$ en el camino lógico. Se valida estrictamente que $u$ sea hijo izquierdo de $v$ y que $u$ carezca de hijo derecho (\texttt{v.left is u} y \texttt{u.right is None}). De no cumplirse ambas condiciones, se rechaza la desconexión con \texttt{ValueError}.

\smallskip
\noindent\textbf{8. ¿Por qué \texttt{PATH-SUM(u, v)} se resuelve tras re-enraizar en $u$ y exponer $v$?} \\
Al designar a $u$ como raíz representada y luego exponer a $v$, el camino simple entre $u$ y $v$ queda confinado con exactitud en el árbol auxiliar cuya raíz es $v$. Gracias al mantenimiento en tiempo real del acumulador en cada nodo, el valor \texttt{v.path\_sum} contiene directamente la suma agregada del camino en tiempo $O(1)$.

\smallskip
\noindent\textbf{9. ¿Qué patrones de entrada generaron más rotaciones por operación?}\\
La \textbf{cadena lineal}, alcanzando un promedio de $23.98$ rotaciones por operación para $n = 1\,600$. En estructuras lineales, los caminos tienen longitud $O(n)$, obligando a enlazar múltiples árboles auxiliares profundos en cada exposición. En contraste, la topología estrella mantuvo un costo cuasi-constante de ${\approx} 7.4$ rot/op debido a su profundidad máxima de 2.

\smallskip
\noindent\textbf{10. ¿Por qué una operación individual puede costar $\Omega(n)$ y aun así la estructura ofrecer $O(\log n)$ amortizado?}\\
Una operación aislada puede costar tiempo lineal $\Omega(n)$ cuando recorre una cadena auxiliar desbalanceada. Sin embargo, las rotaciones zig-zig reducen a la mitad la profundidad de todos los nodos transitados, liberando energía potencial acumulada (según la función potencial de Sleator-Tarjan). Este potencial financia las operaciones costosas y garantiza el promedio $O(\log n)$ sobre cualquier secuencia.

\smallskip
\noindent\textbf{11. ¿Qué errores de implementación detectó el modelo BFS/DFS?}\\
El oráculo BFS/DFS detectó dos fallas sutiles durante la fase de depuración: (1)~\textit{Cortes espurios}: inicialmente \texttt{cut} solo verificaba conectividad, permitiendo cortar pares de nodos conectados por caminos indirectos sin arista física directa; y (2)~\textit{Omisión de actualización}: tras ciertas rotaciones se omitía la llamada a \texttt{\_update}, generando valores de \texttt{path\_sum} inconsistentes.

\smallskip
\noindent\textbf{12. ¿Qué aspectos de esta implementación académica limitarían su uso directo en un sistema real?}\\
Identificamos tres factores limitantes: (1)~\textit{Sobrecarga de punteros en Python}: el almacenamiento en objetos dinámicos dispersos en el montón (\textit{heap}) degrada el rendimiento de la memoria caché frente a implementaciones contiguas en C/C++; (2)~\textit{Límite de recursión}: llamadas recursivas en árboles desbalanceados pueden alcanzar el tope del intérprete (\texttt{sys.getrecursionlimit()}); y (3)~\textit{Falta de concurrencia}: las consultas de solo lectura modifican activamente los punteros por los autoajustes splay, impidiendo accesos concurrentes sin bloqueos de exclusión mutua.
